%! TEX program = xelatex
% MAU LOI GIAI — cac cho dang (hai dau nho hon) ten_khoa (hai dau lon hon) duoc docs/loi_giai/sinh_loi_giai.py thay bang so lieu tinh tu interpolation/*.py
% Khong sua loi_giai_bai_tap.tex truc tiep: sua file nay roi chay lai sinh_loi_giai.py.
\documentclass[11pt,a4paper]{article}
\usepackage[a4paper,margin=1.8cm]{geometry}
\usepackage{amsmath,amssymb,mathtools}
\usepackage{fontspec}
\setmainfont{Times New Roman}
\usepackage{array,booktabs,adjustbox,enumitem,xcolor}
\usepackage[hidelinks]{hyperref}
\usepackage{titlesec}
\usepackage{fancyhdr}

\setlength{\fboxsep}{1.5pt}
\newcommand{\ob}[1]{\fbox{$#1$}}
\newsavebox{\klbox}
\newenvironment{ketluan}
  {\par\medskip\noindent\begin{lrbox}{\klbox}\begin{minipage}{\dimexpr\linewidth-12pt\relax}\textbf{Kết luận. }}
  {\end{minipage}\end{lrbox}\setlength{\fboxsep}{5pt}\fbox{\usebox{\klbox}}\par\medskip}
% Thuat toan: moi buoc mot dong, nhan "Buoc k:" thang cot, dong tran thut vao duoi noi dung
\newenvironment{thuattoan}[1]
  {\par\medskip\noindent\textbf{Thuật toán: #1}\begin{description}[style=standard,align=left,leftmargin=4.3em,labelwidth=3.9em,labelsep=0.4em,itemsep=1pt,topsep=3pt,font=\normalfont\bfseries]}
  {\end{description}}
% Buoc con (2.1, 2.2 / TH1, TH2) trong mot buoc
\newenvironment{con}
  {\begin{description}[style=standard,align=left,leftmargin=2.1em,labelwidth=1.8em,labelsep=0.3em,itemsep=0pt,topsep=1pt,font=\normalfont]}
  {\end{description}}
\newcommand{\pp}[1]{\par\smallskip\noindent\textbf{#1}\ }
% Trong mot buoc: moi y mot dong (\newline); than vong lap / nhanh xuong dong va thut vao (\thut)
\newcommand{\thut}{\newline\hspace*{1.2em}}
\titleformat{\section}{\large\bfseries}{}{0pt}{}
\titleformat{\subsection}{\normalsize\bfseries}{}{0pt}{}
\titlespacing*{\section}{0pt}{14pt}{6pt}
\setlist{itemsep=1pt,topsep=2pt}
\pagestyle{fancy}\fancyhf{}\rhead{Phương pháp số MI3042}\cfoot{\thepage}
\renewcommand{\headrulewidth}{0.4pt}
\allowdisplaybreaks
\setlength{\emergencystretch}{3em}

\begin{document}

\begin{center}
{\Large\bfseries LỜI GIẢI BÀI TẬP PHƯƠNG PHÁP SỐ (MI3042)}\\[2pt]
{\large Câu 3 -- Câu 23 \quad(Chương 1: Xấp xỉ hàm số)}
\end{center}

\noindent\textbf{Quy ước chung.}
\begin{itemize}
\item Đa thức $P_n(x)=a_nx^n+\dots+a_1x+a_0$: chỉ số $=$ số mũ; vector hệ số viết \textbf{từ bậc cao xuống bậc thấp} $(a_n;\dots;a_0)$, vd $3x^3-2x-1\leftrightarrow(3;\,0;\,-2;\,-1)$.
\item Trong các bảng (tỷ) sai phân, \textbf{các ô được chọn để lập đa thức được khoanh khung}, vd \ob{0.1234}. Trong bảng Horner, ô khung là số dư $r_j$.
\item Bảng tích: hàng đầu là đa thức hằng $1$; mỗi hàng sau nhân hàng trên với $(x-\text{mốc ghi ở cột đầu})$ bằng Horner nhân: với $A_{k,l}$ là hệ số của $x^l$ ở hàng $k$ và $z$ là mốc ở cột đầu hàng $k$ thì $A_{0,0}=1$, $A_{k,l}=A_{k-1,l-1}-z\,A_{k-1,l}$. Các hàng hệ số của bảng tích tạo thành ma trận $A$; với $C$ là vector hàng các hệ số $c_k$ thì vector hệ số của $\sum_kc_k\omega_k$ là $C\cdot A$.
\item Số tính ra ghi tối đa 6 chữ số thập phân; số có $|v|<10^{-4}$ ghi dạng $a\cdot10^k$ với 6 chữ số thập phân ở $a$. Số liệu đề cho và sai phân của số liệu cách đều (hiệu các số cùng số chữ số thập phân) ghi đúng, không làm tròn.
\end{itemize}

%=====================================================================
\section{Câu 3. Mốc nội suy tối ưu}

\subsection{a) Xây dựng công thức $n$ mốc nội suy tối ưu trên $[-1,1]$ và trên $[a,b]$}

\pp{Sai số nội suy.} Với $n$ mốc $x_0,\dots,x_{n-1}\in[a,b]$ và $P$ bậc $\le n-1$:
\[ f(x)-P(x)=\frac{f^{(n)}(\xi)}{n!}\,w_n(x),\qquad w_n(x)=\prod_{i=0}^{n-1}(x-x_i). \]
$w_n$ là đa thức bậc $n$ có hệ số cao nhất bằng 1; mốc tối ưu là mốc làm $\max_{[a,b]}|w_n(x)|$ nhỏ nhất.

\pp{Trên $[-1,1]$.} Đa thức Chebyshev $T_n(t)=\cos(n\arccos t)$ có hệ số cao nhất $2^{n-1}$ và $|T_n(t)|\le1$. Trong các đa thức bậc $n$ có hệ số cao nhất bằng 1, $\tilde T_n=2^{1-n}T_n$ có $\max_{[-1,1]}|\tilde T_n|=2^{1-n}$ nhỏ nhất. Vậy chọn $w_n=\tilde T_n$, tức các mốc là nghiệm của $T_n$:
\[ \boxed{t_i=\cos\frac{(2i+1)\pi}{2n},\ i=0,\dots,n-1,\qquad \min\max_{[-1,1]}|w_n|=\frac{1}{2^{n-1}}.} \]

\pp{Trên $[a,b]$.} Đổi biến $x=\dfrac{a+b}{2}+\dfrac{b-a}{2}t$ ánh xạ $[-1,1]$ lên $[a,b]$. Khi đó $x-x_i=\frac{b-a}{2}(t-t_i)$ nên $w_n(x)=\left(\frac{b-a}{2}\right)^n\prod(t-t_i)$; bài toán trên $[a,b]$ tương đương bài toán trên $[-1,1]$:
\[ \boxed{x_i=\frac{a+b}{2}+\frac{b-a}{2}\cos\frac{(2i+1)\pi}{2n},\qquad \min\max_{[a,b]}|w_n|=\left(\frac{b-a}{2}\right)^n\frac{1}{2^{n-1}}=\frac{(b-a)^n}{2^{2n-1}},} \]
và khi đó $|f(x)-P(x)|\le\dfrac{M_n}{n!}\cdot\dfrac{(b-a)^n}{2^{2n-1}}$, $M_n=\max_{[a,b]}|f^{(n)}|$.

\pp{Sắp tăng dần.} $\cos$ giảm trên $[0,\pi]$ nên $t_i$ ở trên giảm theo $i$. Dùng $t_i=-\cos\frac{(2i+1)\pi}{2n}=\cos\frac{(2(n-1-i)+1)\pi}{2n}$ (cùng tập nghiệm, đánh số ngược) thì $t_0<t_1<\dots<t_{n-1}$.

\subsection{b) Thuật toán xác định $n$ mốc tối ưu sắp tăng dần trên $[a,b]$}
\begin{thuattoan}{Mốc nội suy tối ưu}
\item[Input:] số mốc $n$; đoạn $[a,b]$.
\item[Output:] mảng $x=(x_0,\dots,x_{n-1})$, $x_0<x_1<\dots<x_{n-1}$.
\item[Bước 0:] Nếu $n$ không nguyên hoặc $n<1$ thì báo lỗi, dừng thuật toán.\newline Nếu $a=b$ thì báo lỗi, dừng thuật toán.\newline Nếu $a>b$ thì đổi chỗ $a$ và $b$.
\item[Bước 1:] Lặp $i=0,\dots,n-1$: tính $x_i:=\dfrac{a+b}{2}-\dfrac{b-a}{2}\cos\dfrac{(2i+1)\pi}{2n}$.
\item[Bước 2:] Dừng thuật toán, trả về $x$.
\end{thuattoan}

\subsection{c) Áp dụng}
\pp{Công thức:} $x_i=\dfrac{a+b}{2}-\dfrac{b-a}{2}\cos\dfrac{(2i+1)\pi}{2n}$, $i=0,\dots,n-1$.

\pp{$n=9$, $[a,b]=[-1,1]$:} $x_i=-\cos\dfrac{(2i+1)\pi}{18}$, $i=0,\dots,8$:
<<c3_bang_a>>
<<c3_x_a>>

\pp{$n=12$, $[a,b]=[3,6]$:} $x_i=4.5-1.5\cos\dfrac{(2i+1)\pi}{24}$, $i=0,\dots,11$:
<<c3_bang_b>>
<<c3_x_b>>

%=====================================================================
\section{Câu 4. Sơ đồ Horner: giá trị, phép chia, đạo hàm}

\subsection{a) Giá trị $P_n(c)$ \quad b) Chia $P_n(x)$ cho $(x-c)$}
\begin{thuattoan}{Horner chia (giá trị, thương, số dư)}
\item[Input:] vector hệ số $a=(a_n,\dots,a_0)$; số $c$.
\item[Output:] thương $b=(b_{n-1},\dots,b_0)$; số dư $r=P_n(c)$.
\item[Bước 0:] Bỏ các hệ số $0$ ở đầu vector $a$.\newline Nếu không còn hệ số nào thì báo $P_n\equiv0$, dừng thuật toán.\newline Đặt $n:=(\text{số phần tử của }a)-1$.
\item[Bước 1:] Nếu $n=0$ thì đặt $r:=a_0$, dừng thuật toán, trả về $r$.
\item[Bước 2:] Đặt $b_{n-1}:=a_n$.
\item[Bước 3:] Lặp $i=n-1,\dots,1$: tính $b_{i-1}:=a_i+c\,b_i$.
\item[Bước 4:] Tính $r:=a_0+c\,b_0$.\newline Dừng thuật toán, trả về $b$ và $r$.
\end{thuattoan}

\subsection{c) Đạo hàm cấp $k$ của $P_n$ tại $c$}
\begin{thuattoan}{Đạo hàm cấp $0..k$ tại $c$ (Horner chia lặp)}
\item[Input:] vector hệ số $a=(a_n,\dots,a_0)$; số $c$; cấp $k$.
\item[Output:] vector $r=(r_0,\dots,r_K)$, $K=\min(k,n)$; vector $v=(v_0,\dots,v_k)$ với $v_j=P_n^{(j)}(c)$.
\item[Bước 0:] Bỏ các hệ số $0$ ở đầu vector $a$.\newline Nếu không còn hệ số nào thì báo $P_n\equiv0$, dừng thuật toán.\newline Đặt $n:=(\text{số phần tử của }a)-1$.\newline Nếu $k$ không nguyên hoặc $k<0$ thì báo lỗi, dừng thuật toán.
\item[Bước 1:] Đặt $K:=\min(k,n)$.\newline Lặp $i=0,\dots,n$: đặt $b_i:=a_i$.
\item[Bước 2:] Lặp $j=0,\dots,K$:
  \begin{con}
  \item[2.1:] Lặp $i=n-1,\dots,j$: tính $b_i:=b_i+c\,b_{i+1}$ (chia cho $x-c$).
  \item[2.2:] Đặt $r_j:=b_j$. Tính $v_j:=j!\,r_j$.
  \end{con}
\item[Bước 3:] Lặp $j=K+1,\dots,k$: đặt $v_j:=0$.\newline Dừng thuật toán, trả về $r$ và $v$.
\end{thuattoan}

\subsection{d) Áp dụng: $P(x)=3x^3-2x-1$, $a=(3;\,0;\,-2;\,-1)$}
\pp{Công thức} (Horner chia lặp, câu c). Hàng ``Hệ số'' là $b^{(-1)}_i=a_i$; hàng $j$ ghi $b^{(j)}_i$ ($i=3,\dots,j$) sau lần chia thứ $j$, hàng $c$ phía trên ghi $c\,b^{(j)}_{i+1}$:
\[ b^{(j)}_3=a_3,\qquad b^{(j)}_i=b^{(j-1)}_i+c\,b^{(j)}_{i+1}\ \ (i=2,\dots,j),\qquad r_j=b^{(j)}_j,\qquad P^{(j)}(c)=j!\,r_j,\qquad j=0,\dots,3. \]
Lần chia $j=0$ cho thương $b^{(0)}_3x^2+b^{(0)}_2x+b^{(0)}_1$ và số dư $r_0=P(c)$; $P^{(j)}(c)=0$ với $j\ge4$.

\pp{$c=0$:}
<<c4_bang_0>>
$P(0)=r_0=<<c4_P_0>>$; thương $<<c4_Q_0>>$, dư $<<c4_P_0>>$;\\
$(P,P',P'',P''')(0)=(0!\,r_0;\ 1!\,r_1;\ 2!\,r_2;\ 3!\,r_3)=<<c4_v_0>>$.

\pp{$c=2.15$:}
<<c4_bang_1>>
$P(2.15)=r_0=<<c4_P_1>>$; thương $<<c4_Q_1>>$, dư $<<c4_P_1>>$;\\
$(P,P',P'',P''')(2.15)=(0!\,r_0;\ 1!\,r_1;\ 2!\,r_2;\ 3!\,r_3)=<<c4_v_1>>$.

%=====================================================================
\section{Câu 5. Phép nhân Horner, đa thức $w_{n+1}(x)$}

\subsection{a) Nhân $P_n(x)$ với $(x-c)$}
\begin{thuattoan}{Horner nhân}
\item[Input:] vector hệ số $a=(a_n,\dots,a_0)$; số $c$.
\item[Output:] vector $d=(d_{n+1},\dots,d_0)$ với $P_n(x)\,(x-c)=\sum_{i=0}^{n+1}d_ix^i$.
\item[Bước 0:] Bỏ các hệ số $0$ ở đầu vector $a$.\newline Nếu không còn hệ số nào thì báo $P_n\equiv0$, dừng thuật toán.\newline Đặt $n:=(\text{số phần tử của }a)-1$.
\item[Bước 1:] Đặt $d_{n+1}:=a_n$.
\item[Bước 2:] Lặp $i=n,\dots,1$: tính $d_i:=a_{i-1}-c\,a_i$.
\item[Bước 3:] Tính $d_0:=-c\,a_0$.\newline Dừng thuật toán, trả về $d$.
\end{thuattoan}

\subsection{b) Hệ số của $w_{n+1}(x)=\prod_{k=0}^{n}(x-x_k)$}
\begin{thuattoan}{Bảng tích (nhân Horner lặp)}
\item[Input:] vector mốc $x=(x_0,\dots,x_n)$.
\item[Output:] vector $p=(p_{n+1},\dots,p_0)$, $p_{n+1}=1$, với $w_{n+1}(x)=\sum_{l=0}^{n+1}p_lx^l$.
\item[Bước 0:] Nếu không có mốc nào thì báo lỗi, dừng thuật toán.\newline Đặt $n:=(\text{số mốc})-1$.
\item[Bước 1:] Đặt $p_0:=1$.
\item[Bước 2:] Lặp $i=0,\dots,n$:
  \begin{con}
  \item[2.1:] Đặt $p_{i+1}:=p_i$.
  \item[2.2:] Lặp $l=i,\dots,1$: tính $p_l:=p_{l-1}-x_i\,p_l$.
  \item[2.3:] Tính $p_0:=-x_i\,p_0$.
  \end{con}
\item[Bước 3:] Dừng thuật toán, trả về $p$.
\end{thuattoan}

\subsection{c) Áp dụng: $x_k=2;\ 2.4;\ 2.7;\ 3;\ 3.1;\ 3.4$}
\pp{Công thức} (bảng tích, câu b): $A_{k,l}$ là hệ số của $x^l$ ở hàng $k$ ($k=0,\dots,6$); hàng $0$ là đa thức $1$, hàng $k$ là hàng $k-1$ nhân với $(x-x_{k-1})$ (cột đầu ghi $x_{k-1}$):
\[ A_{0,0}=1,\qquad A_{k,l}=A_{k-1,l-1}-x_{k-1}\,A_{k-1,l}\quad(l=0,\dots,k;\ A_{k-1,-1}=A_{k-1,k}=0); \]
hàng cuối là $w_6$.
<<c5_bang>>
\begin{ketluan}
Dạng chính tắc của $w_6(x)=(x-2)(x-2.4)(x-2.7)(x-3)(x-3.1)(x-3.4)$:
<<c5_w>>
\end{ketluan}

%=====================================================================
\section{Câu 6. Thuật toán đa thức nội suy Lagrange}

\begin{thuattoan}{Đa thức nội suy Lagrange dạng chính tắc}
\item[Input:] vector $x=(x_0,\dots,x_n)$, $y=(y_0,\dots,y_n)$.
\item[Output:] vector $a=(a_n,\dots,a_0)$ — hệ số của đa thức nội suy Lagrange $P_n(x)=\sum_{l=0}^{n}a_lx^l$.
\item[Bước 0:] Nếu $x$, $y$ khác số phần tử thì báo lỗi, dừng thuật toán.\newline Sắp các mốc tăng theo $x$.\newline Nếu hai mốc trùng nhau thì: cùng $y$ giữ một mốc, khác $y$ báo lỗi, dừng thuật toán.\newline Đặt $n:=(\text{số mốc})-1$.
\item[Bước 1:] Tính vector $p=(p_{n+1},\dots,p_0)$ — hệ số của $w_{n+1}(x)$ — theo thuật toán Câu 5b.
\item[Bước 2:] Khởi tạo vector $a:=(0,\dots,0)$ ($n+1$ phần tử).
\item[Bước 3:] Lặp $i=0,\dots,n$:
  \begin{con}
  \item[3.1:] Chia $w_{n+1}$ cho $(x-x_i)$ theo thuật toán Câu 4a (vector $p$, $c:=x_i$), được vector thương \mbox{$q=(q_n,\dots,q_0)$}.
  \item[3.2:] Đặt $D_i:=1$.\newline Lặp $j=0,\dots,n$, $j\ne i$: tính $D_i:=D_i\,(x_i-x_j)$.
  \item[3.3:] Lặp $l=0,\dots,n$: tính $a_l:=a_l+\dfrac{y_i}{D_i}\,q_l$.
  \end{con}
\item[Bước 4:] Dừng thuật toán, trả về $a$.
\end{thuattoan}

%=====================================================================
\section{Câu 7. Đa thức nội suy Lagrange của bảng số liệu}

\pp{Công thức} (thuật toán Câu 6):
\[ P_5(x)=\sum_{i=0}^{5}\frac{y_i}{D_i}\,q^{(i)}(x),\qquad q^{(i)}(x)=\frac{w_6(x)}{x-x_i},\qquad D_i=\prod_{j\ne i}(x_i-x_j). \]

\pp{Bảng tích} (hàng cuối là hệ số của $w_6=\sum p_lx^l$): $A_{0,0}=1$, $A_{k,l}=A_{k-1,l-1}-x_{k-1}A_{k-1,l}$):
<<c7_bang_tich>>
\pp{Bảng thương} (hàng $i$: chia $w_6$ cho $(x-x_i)$ bằng Horner; $Q_{i,l}$ là hệ số của $x^l$ trong $q^{(i)}$, $p_l$ lấy ở hàng cuối bảng tích):
\[ Q_{i,5}=p_6=1,\qquad Q_{i,l-1}=p_l+x_i\,Q_{i,l}\ \ (l=5,\dots,1),\qquad \text{dư}=p_0+x_i\,Q_{i,0}=0. \]
<<c7_bang_thuong>>
\pp{Bảng hệ số} (ô hàng $i$, cột $j$: $x_i-x_j$; $D_i$ là tích các ô $j\ne i$ của hàng $i$):
<<c7_bang_he_so>>
\pp{Hệ số của $P_5$:} với vector hàng $Y=\bigl(\frac{y_0}{D_0},\dots,\frac{y_5}{D_5}\bigr)$ và ma trận $Q$ (6 cột hệ số của bảng thương), vector hệ số của $P_5$ là $Y\cdot Q$:
<<c7_YQ>>
\begin{ketluan}
<<c7_P>>
\end{ketluan}

%=====================================================================
\section{Câu 8. Đa thức nội suy Newton tiến (mốc bất kỳ)}

\pp{Phương pháp.} Mốc không cách đều nên dùng tỷ sai phân.

\subsection{a) Đa thức Newton tiến}
\pp{Bảng tỷ sai phân} (hàng $i$, cột TSP$k$ là $f[x_{i-k},\dots,x_i]$; đường chéo khung $=c_k$):
\[ f[x_{i-k},\dots,x_i]=\frac{f[x_{i-k+1},\dots,x_i]-f[x_{i-k},\dots,x_{i-1}]}{x_i-x_{i-k}}. \]
<<c8_bang_tsp>>
\pp{Hệ số:}
<<c8_C>>
Dạng Newton tiến $P_5(x)=\sum_{k=0}^{5}c_k\,\omega_k(x)$, $\omega_k(x)=\prod_{i=0}^{k-1}(x-x_i)$:
<<c8_dang_newton>>
\pp{Rút gọn về dạng chính tắc bằng bảng tích} (hàng $k$: $\omega_k=\omega_{k-1}\cdot(x-x_{k-1})$; $A_{k,l}$ là hệ số của $x^l$ trong $\omega_k$):
\[ A_{0,0}=1,\qquad A_{k,l}=A_{k-1,l-1}-x_{k-1}\,A_{k-1,l}\quad(l=0,\dots,k). \]
<<c8_bang_tich>>
Vector hệ số của $P_5(x)$:
<<c8_CA>>
<<c8_P>>

\subsection{b) Tính gần đúng $f(1.8)$, $f'''(1.8)$}
Horner chia lặp $P_5$ tại $c=1.8$ (Câu 4c):
\[ f(1.8)\approx P_5(1.8)=r_0=<<c8_f>>,\qquad f'''(1.8)\approx P_5'''(1.8)=3!\,r_3=6\cdot<<c8_rk>>=<<c8_fk>>. \]

%=====================================================================
\section{Câu 9. Đa thức nội suy Newton lùi (mốc bất kỳ)}

\pp{Phương pháp.} Mốc không cách đều nên dùng tỷ sai phân; Newton lùi xuất phát từ $x_5=3.2$.

\subsection{a) Đa thức Newton lùi}
\pp{Bảng tỷ sai phân} (hàng $i$, cột TSP$k$ là $f[x_{i-k},\dots,x_i]$; hàng cuối khung $=c_k$):
\[ f[x_{i-k},\dots,x_i]=\frac{f[x_{i-k+1},\dots,x_i]-f[x_{i-k},\dots,x_{i-1}]}{x_i-x_{i-k}}. \]
<<c9_bang_tsp>>
\pp{Hệ số:}
<<c9_C>>
Dạng Newton lùi $P_5(x)=\sum_{k=0}^{5}c_k\,\omega_k(x)$, $\omega_k(x)=\prod_{i=0}^{k-1}(x-x_{5-i})$:
<<c9_dang_newton>>
\pp{Bảng tích} (hàng $k$: $\omega_k=\omega_{k-1}\cdot(x-x_{6-k})$; $A_{k,l}$ là hệ số của $x^l$ trong $\omega_k$):
\[ A_{0,0}=1,\qquad A_{k,l}=A_{k-1,l-1}-x_{6-k}\,A_{k-1,l}\quad(l=0,\dots,k). \]
<<c9_bang_tich>>
Vector hệ số của $P_5(x)$:
<<c9_CA>>
<<c9_P>>

\subsection{b) Tính gần đúng $f(2.4)$, $f^{(4)}(2.4)$}
Horner chia lặp $P_5$ tại $c=2.4$ (Câu 4c):
\[ f(2.4)\approx P_5(2.4)=r_0=<<c9_f>>,\qquad f^{(4)}(2.4)\approx P_5^{(4)}(2.4)=4!\,r_4=24\cdot<<c9_rk>>=<<c9_fk>>. \]

%=====================================================================
\section{Câu 10. Thuật toán đa thức nội suy Newton tiến / lùi}

\begin{thuattoan}{Newton tiến / lùi, mốc bất kỳ}
\item[Input:] vector $x=(x_0,\dots,x_n)$, $y=(y_0,\dots,y_n)$; chiều (tiến hoặc lùi); điểm $\bar x$; cấp đạo hàm $r$.
\item[Output:] vector $a=(a_n,\dots,a_0)$ với $P_n(x)=\sum_la_lx^l$; vector $f=(f_0,\dots,f_r)$, $f_j\approx f^{(j)}(\bar x)$.
\item[Bước 0:] Nếu $x$, $y$ khác số phần tử thì báo lỗi, dừng thuật toán.\newline Sắp các mốc tăng theo $x$.\newline Nếu hai mốc trùng nhau thì: cùng $y$ giữ một mốc, khác $y$ báo lỗi, dừng thuật toán.\newline Đặt $n:=(\text{số mốc})-1$.\newline Nếu $r$ không nguyên hoặc $r<0$ thì báo lỗi, dừng thuật toán.
\item[Bước 1:] Khởi tạo ma trận $T$ cỡ $(n+1)\times(n+1)$ ($T_{i,k}=f[x_{i-k},\dots,x_i]$).\newline Lặp $i=0,\dots,n$: đặt $T_{i,0}:=y_i$.\newline Lặp lồng $k=1,\dots,n$; $i=k,\dots,n$: tính $T_{i,k}:=\dfrac{T_{i,k-1}-T_{i-1,k-1}}{x_i-x_{i-k}}$.
\item[Bước 2:] Lặp $k=0,\dots,n$:\thut nếu tiến thì đặt $c_k:=T_{k,k}$;\thut nếu lùi thì đặt $c_k:=T_{n,k}$.
\item[Bước 3:] Khởi tạo vector $e$, $a$, mỗi vector $n+1$ phần tử (chỉ số $=$ số mũ):\thut $(e_0,e_1,\dots,e_n):=(1,0,\dots,0)$;\thut $(a_0,a_1,\dots,a_n):=(c_0,0,\dots,0)$.\newline Lặp $k=1,\dots,n$:
  \begin{con}
  \item[3.1:] Nếu tiến thì đặt $z:=x_{k-1}$; nếu lùi thì đặt $z:=x_{n-k+1}$.\newline Nhân vector $e$ với $(x-z)$ theo thuật toán Câu 5a ($e$ thành hệ số của $\omega_k(x)$).
  \item[3.2:] Lặp $l=0,\dots,k$: tính $a_l:=a_l+c_k\,e_l$.
  \end{con}
\item[Bước 4:] Gọi thuật toán Câu 4c với vector $a$, $c:=\bar x$, cấp $k:=r$, được vector $f=(f_0,\dots,f_r)$, $f_j=P_n^{(j)}(\bar x)$.\newline Dừng thuật toán, trả về $a$ và $f$.
\end{thuattoan}

\begin{thuattoan}{Newton tiến / lùi, mốc cách đều}
\item[Input:] vector $x=(x_0,\dots,x_n)$ cách đều, $y=(y_0,\dots,y_n)$; chiều (tiến hoặc lùi); điểm $\bar x$; cấp đạo hàm $r$.
\item[Output:] vector $b=(b_n,\dots,b_0)$ với $P_n(t)=\sum_lb_lt^l$ ($t=\frac{x-x_0}{h}$ nếu tiến, $t=\frac{x-x_n}{h}$ nếu lùi); vector $f=(f_0,\dots,f_r)$, $f_j\approx f^{(j)}(\bar x)$.
\item[Bước 0:] Nếu $x$, $y$ khác số phần tử thì báo lỗi, dừng thuật toán.\newline Sắp các mốc tăng theo $x$.\newline Nếu hai mốc trùng nhau thì: cùng $y$ giữ một mốc, khác $y$ báo lỗi, dừng thuật toán.\newline Đặt $n:=(\text{số mốc})-1$.\newline Nếu $n<1$ thì báo lỗi (cần ít nhất 2 mốc), dừng thuật toán.\newline Nếu $r$ không nguyên hoặc $r<0$ thì báo lỗi, dừng thuật toán.
\item[Bước 1:] Đặt $h:=\dfrac{x_n-x_0}{n}$.\newline Lặp $i=0,\dots,n-1$:\thut nếu $|x_{i+1}-x_i-h|>10^{-6}h$ thì báo ``mốc không cách đều'', dùng thuật toán mốc bất kỳ, dừng thuật toán.
\item[Bước 2:] Khởi tạo ma trận $T$ cỡ $(n+1)\times(n+1)$ ($T_{i,k}=\Delta^ky_{i-k}$).\newline Lặp $i=0,\dots,n$: đặt $T_{i,0}:=y_i$.\newline Lặp lồng $k=1,\dots,n$; $i=k,\dots,n$: tính $T_{i,k}:=T_{i,k-1}-T_{i-1,k-1}$.
\item[Bước 3:] Lặp $k=0,\dots,n$:\thut nếu tiến thì tính $c_k:=T_{k,k}/k!$;\thut nếu lùi thì tính $c_k:=T_{n,k}/k!$.
\item[Bước 4:] Khởi tạo vector $e$, $b$, mỗi vector $n+1$ phần tử (chỉ số $=$ số mũ của $t$):\thut $(e_0,e_1,\dots,e_n):=(1,0,\dots,0)$;\thut $(b_0,b_1,\dots,b_n):=(c_0,0,\dots,0)$.\newline Lặp $k=1,\dots,n$:
  \begin{con}
  \item[4.1:] Nếu tiến thì đặt $z:=k-1$; nếu lùi thì đặt $z:=-(k-1)$.\newline Nhân vector $e$ với $(t-z)$ theo thuật toán Câu 5a ($e$ thành hệ số của $\omega_k(t)$).
  \item[4.2:] Lặp $l=0,\dots,k$: tính $b_l:=b_l+c_k\,e_l$.
  \end{con}
\item[Bước 5:] Nếu tiến thì tính $\bar t:=(\bar x-x_0)/h$; nếu lùi thì tính $\bar t:=(\bar x-x_n)/h$.
\item[Bước 6:] Gọi thuật toán Câu 4c với vector $b$, $c:=\bar t$, cấp $k:=r$, được vector $d=(d_0,\dots,d_r)$.\newline Lặp $j=0,\dots,r$: tính $f_j:=d_j/h^j$.\newline Dừng thuật toán, trả về $b$ và $f$.
\end{thuattoan}

%=====================================================================
\section{Câu 11. Trích xuất mốc và nội suy Newton trên bảng dài}

\subsection{a) Thuật toán trích xuất $k$ điểm nội suy phù hợp để tính $f(\bar x)$}
\begin{thuattoan}{Trích xuất $k$ mốc gần $\bar x$ nhất}
\item[Input:] vector $x=(x_0,\dots,x_n)$ cách đều, $y=(y_0,\dots,y_n)$; điểm $\bar x$; số mốc $k$.
\item[Output:] chỉ số $L$, $R$ ($R-L+1=k$); vector $X=(x_L,\dots,x_R)$, $Y=(y_L,\dots,y_R)$.
\item[Bước 0:] Nếu $x$, $y$ khác số phần tử thì báo lỗi, dừng thuật toán.\newline Sắp các mốc tăng theo $x$.\newline Nếu hai mốc trùng nhau thì: cùng $y$ giữ một mốc, khác $y$ báo lỗi, dừng thuật toán.\newline Đặt $n:=(\text{số mốc})-1$.\newline Đặt $h:=\dfrac{x_n-x_0}{n}$.\newline Lặp $i=0,\dots,n-1$:\thut nếu $|x_{i+1}-x_i-h|>10^{-6}h$ thì báo ``mốc không cách đều'', dừng thuật toán.\newline Nếu $k$ không nguyên hoặc $k<1$ thì báo lỗi, dừng thuật toán.\newline Nếu $k>n+1$ thì báo ``không đủ mốc'', đặt $k:=n+1$.\newline Nếu $\bar x<x_0$ hoặc $\bar x>x_n$ thì cảnh báo ngoại suy.
\item[Bước 1:] Tìm mốc gần $\bar x$ nhất: tính $j:=\operatorname{round}\dfrac{\bar x-x_0}{h}$ (số nguyên gần nhất).\newline Nếu $j<0$ thì đặt $j:=0$; nếu $j>n$ thì đặt $j:=n$.\newline Đặt $L:=j$, $R:=j$ (các mốc đã chọn là $x_L,\dots,x_R$).
\item[Bước 2:] Lặp $s=2,\dots,k$ (mỗi lần thêm mốc kề bên trái hoặc phải, phía gần $\bar x$ hơn):\thut nếu $L=0$ thì $R:=R+1$;\thut ngược lại nếu $R=n$ thì $L:=L-1$;\thut ngược lại nếu $\bar x-x_{L-1}\le x_{R+1}-\bar x$ thì $L:=L-1$, còn không thì $R:=R+1$.
\item[Bước 3:] Đặt $X:=(x_L,\dots,x_R)$, $Y:=(y_L,\dots,y_R)$.\newline Dừng thuật toán, trả về $L$, $R$, $X$, $Y$.
\end{thuattoan}

\subsection{b) Trích xuất 9 điểm cho $\bar x=3.52$}
\pp{Xử lý dữ liệu.} File \texttt{data Câu 11.xlsx} (sheet đầu): cột A là $x$, cột B là $y$; ô B1 $=2.95$ không có $x$ đi kèm nên bị loại (dòng 1, ô $x$ rỗng). Còn $<<c11_N>>$ mốc $x=<<c11_x0>>,\dots,<<c11_xn>>$, cách đều $h=<<c11_h>>$.

Thuật toán câu a ($\bar x=3.52$, $k=9$): <<c11_ket_trich>> Đánh số lại $i=0..8$:
<<c11_bang_xy>>

\subsection{c) Đa thức Newton tiến, $f(\bar x)$, $f'(\bar x)$}
\pp{Phương pháp.} 9 mốc cách đều $h=0.1$ nên dùng Newton tiến mốc cách đều, $t=\dfrac{x-x_0}{h}=\dfrac{x-3.1}{0.1}$.

\pp{Bảng sai phân} (hàng $i$, cột $\Delta^k$ là $\Delta^ky_{i-k}$; đường chéo khung $=\Delta^ky_0$):
\[ \Delta^0y_i=y_i,\qquad \Delta^ky_i=\Delta^{k-1}y_{i+1}-\Delta^{k-1}y_i\quad(k=1,\dots,8;\ i=0,\dots,8-k). \]
<<c11_bang_sp>>
\pp{Hệ số:}
<<c11_C>>
\pp{Bảng tích theo $t$} (hàng $k$: $\omega_k(t)=\omega_{k-1}(t)\cdot(t-z_{k-1})$, $z_{k-1}=k-1$, tức $\omega_k(t)=t(t-1)\cdots(t-k+1)$; $A_{k,l}$ là hệ số của $t^l$):
\[ A_{0,0}=1,\qquad A_{k,l}=A_{k-1,l-1}-(k-1)\,A_{k-1,l}\quad(l=0,\dots,k). \]
<<c11_bang_tich>>
Vector hệ số của $P_8(t)=\sum_{k=0}^{8}c_k\,\omega_k(t)$ là $C\cdot A$:
<<c11_CA>>
<<c11_Pt>>

\pp{Tính tại $\bar x=3.52$:} $\bar t=\dfrac{\bar x-x_0}{h}=<<c11_t>>$. Horner chia lặp $P_8$ tại $\bar t$ (Câu 4c):
\[ f(3.52)\approx P_8(\bar t)=<<c11_f>>,\qquad f'(3.52)\approx\frac{P_8'(\bar t)}{h}=\frac{<<c11_dP>>}{0.1}=<<c11_df>>. \]

%=====================================================================
\section{Câu 12. Đa thức nội suy Gauss I}

\pp{Phương pháp.} 9 mốc cách đều, $\bar t=0.8>0$ nên dùng Gauss I, mốc trung tâm $x_0=<<c12_x0>>$ ($i=<<c12_p>>..<<c12_pn>>$), $h=0.1$, $t=\frac{x-x_0}{h}$.

\subsection{a) Lập đa thức Gauss I}
\[ P(x_0+th)=y_0+t\Delta y_0+\frac{\Delta^2y_{-1}}{2!}t(t-1)+\frac{\Delta^3y_{-1}}{3!}(t+1)t(t-1)+\frac{\Delta^4y_{-2}}{4!}(t+1)t(t-1)(t-2)+\cdots \]
\pp{Bảng sai phân} (hàng $i$ chứa $y_i,\Delta y_i,\Delta^2y_i,\dots$; ô khung: $\Delta^ky_{-\lfloor k/2\rfloor}$):
\[ \Delta^0y_i=y_i,\qquad \Delta^ky_i=\Delta^{k-1}y_{i+1}-\Delta^{k-1}y_i\quad(k=1,\dots,8;\ i=-4,\dots,4-k). \]
<<c12_bang_sp>>
\pp{Hệ số:}
<<c12_C>>
\pp{Bảng tích} (hàng $k$: $\omega_k(t)=\omega_{k-1}(t)\cdot(t-z_{k-1})$, $z=0,1,-1,2,-2,3,-3,4$; $A_{k,l}$ là hệ số của $t^l$):
\[ A_{0,0}=1,\qquad A_{k,l}=A_{k-1,l-1}-z_{k-1}\,A_{k-1,l}\quad(l=0,\dots,k). \]
<<c12_bang_tich>>
Vector hệ số của $P_8(t)=\sum_{k=0}^{8}c_k\,\omega_k(t)$ là $C\cdot A$:
<<c12_CA>>
<<c12_Pt>>
với $t=(x-9.6)/0.1$.

\subsection{b) Tính gần đúng $f(9.68)$}
$\bar t=\dfrac{\bar x-x_0}{h}=<<c12_t>>$. Horner chia $P_8$ tại $\bar t$ (Câu 4a): $f(9.68)\approx P_8(\bar t)=<<c12_f>>$.

%=====================================================================
\section{Câu 13. Đa thức nội suy Stirling}

\pp{Phương pháp.} 9 mốc cách đều (số lẻ) nên dùng Stirling, mốc trung tâm $x_0=<<c13_x0>>$ ($i=<<c13_p>>..<<c13_pn>>$), $h=0.1$, $t=\frac{x-x_0}{h}$.

\subsection{a) Lập đa thức Stirling}
\[ P(x_0+th)=y_0+t\,\frac{\Delta y_{-1}+\Delta y_0}{2}+\frac{t^2}{2!}\Delta^2y_{-1}+\frac{t(t^2-1)}{3!}\cdot\frac{\Delta^3y_{-2}+\Delta^3y_{-1}}{2}+\frac{t^2(t^2-1)}{4!}\Delta^4y_{-2}+\cdots \]
\pp{Bảng sai phân} (ô khung: cặp $\Delta^{2j-1}y_{-j},\Delta^{2j-1}y_{-j+1}$ và $\Delta^{2j}y_{-j}$):
\[ \Delta^0y_i=y_i,\qquad \Delta^ky_i=\Delta^{k-1}y_{i+1}-\Delta^{k-1}y_i\quad(k=1,\dots,8;\ i=-4,\dots,4-k). \]
<<c13_bang_sp>>
\pp{Hệ số:}
<<c13_C>>
\pp{Bảng tích theo $u=t^2$} (hàng $j$: $\Omega_j(u)=\Omega_{j-1}(u)\cdot(u-j^2)$; $A_{j,l}$ là hệ số của $u^l$):
\[ A_{0,0}=1,\qquad A_{j,l}=A_{j-1,l-1}-j^2A_{j-1,l}\quad(l=0,\dots,j). \]
<<c13_bang_tich>>
<<c13_CA>>
Ghép hệ số (xen kẽ chẵn, lẻ) theo $P(t)=y_0+t\,P_{\text{lẻ}}(t^2)+t^2P_{\text{chẵn}}(t^2)$:
<<c13_Pt>>
với $t=(x-1.8)/0.1$.

\subsection{b) Giá trị và đạo hàm tại $\bar x=1.63$}
$\bar t=\dfrac{\bar x-x_0}{h}=<<c13_t>>$. Horner chia lặp $P_8$ tại $\bar t$ (Câu 4c):
\[ f(1.63)\approx P_8(\bar t)=<<c13_f>>,\qquad f'(1.63)\approx\frac{P_8'(\bar t)}{h}=\frac{<<c13_dP>>}{0.1}=<<c13_df>>. \]

%=====================================================================
\section{Câu 14. Đa thức nội suy Bessel}

\pp{Đề (đã sửa).} Câu 14b ghi $\bar x=1.63$ (ngoài bảng $[5;5.7]$); sửa thành $\bar x=5.23$.

\pp{Phương pháp.} 8 mốc cách đều (số chẵn) nên dùng Bessel, hai mốc chính giữa $x_0=<<c14_x0>>$, $x_1=5.4$ ($i=<<c14_p>>..<<c14_pn>>$), $h=0.1$, $t=\frac{x-x_0}{h}$.

\subsection{a) Lập đa thức Bessel}
\begin{align*}
P(x_0+th)&=\frac{y_0+y_1}{2}+\Big(t-\frac12\Big)\Delta y_0+\frac{t(t-1)}{2!}\cdot\frac{\Delta^2y_{-1}+\Delta^2y_0}{2}+\frac{(t-\frac12)t(t-1)}{3!}\Delta^3y_{-1}\\
&\quad+\frac{(t+1)t(t-1)(t-2)}{4!}\cdot\frac{\Delta^4y_{-2}+\Delta^4y_{-1}}{2}+\cdots
\end{align*}
\pp{Bảng sai phân} (ô khung: cặp $\Delta^{2j}y_{-j},\Delta^{2j}y_{-j+1}$ và $\Delta^{2j+1}y_{-j}$):
\[ \Delta^0y_i=y_i,\qquad \Delta^ky_i=\Delta^{k-1}y_{i+1}-\Delta^{k-1}y_i\quad(k=1,\dots,7;\ i=-3,\dots,4-k). \]
<<c14_bang_sp>>
\pp{Hệ số:}
<<c14_C>>
\pp{Bảng tích theo $u=v^2$, $v=t-\frac12$} (hàng $j$: $\Omega_j(u)=\Omega_{j-1}(u)\cdot\big(u-(j-\frac12)^2\big)$; $A_{j,l}$ là hệ số của $u^l$):
\[ A_{0,0}=1,\qquad A_{j,l}=A_{j-1,l-1}-\big(j-\tfrac12\big)^2A_{j-1,l}\quad(l=0,\dots,j). \]
<<c14_bang_tich>>
<<c14_CA>>
Ghép hệ số (xen kẽ lẻ, chẵn) theo $P(v)=P_{\text{chẵn}}(v^2)+v\,P_{\text{lẻ}}(v^2)$:
<<c14_Pt>>
với $v=t-\frac12=(x-<<c14_s0>>)/0.1$.

\subsection{b) Giá trị và đạo hàm tại $\bar x=5.23$}
$\bar t=\dfrac{\bar x-x_0}{h}=<<c14_t>>$, $\bar v=\bar t-\frac12=<<c14_v>>$. Horner chia lặp $P_7$ tại $\bar v$ (Câu 4c):
\[ f(5.23)\approx P_7(\bar v)=<<c14_f>>,\qquad f'(5.23)\approx\frac{P_7'(\bar v)}{h}=\frac{<<c14_dP>>}{0.1}=<<c14_df>>. \]

%=====================================================================
\section{Câu 15. Thuật toán Gauss I, Gauss II, Stirling, Bessel}

\subsection{a) Gauss I, Gauss II}
\begin{thuattoan}{Gauss I / Gauss II}
\item[Input:] vector $x$, $y$ (mốc cách đều); điểm $\bar x$; cấp đạo hàm $r$; số mốc dùng $m$; công thức (Gauss I hoặc Gauss II).
\item[Output:] vector $b=(b_n,\dots,b_0)$ với $P(t)=\sum_lb_lt^l$, $t=\frac{x-x_0}{h}$; vector $f=(f_0,\dots,f_r)$, $f_j\approx f^{(j)}(\bar x)$.
\item[Bước 0:] Nếu $x$, $y$ khác số phần tử thì báo lỗi, dừng thuật toán.\newline Sắp các mốc tăng theo $x$.\newline Nếu hai mốc trùng nhau thì: cùng $y$ giữ một mốc, khác $y$ báo lỗi, dừng thuật toán.\newline Đặt $n:=(\text{số mốc})-1$.\newline Đặt $h:=\dfrac{x_n-x_0}{n}$.\newline Lặp $i=0,\dots,n-1$:\thut nếu $|x_{i+1}-x_i-h|>10^{-6}h$ thì báo ``mốc không cách đều'', dừng thuật toán.\newline Nếu $m$, $r$ không nguyên, $m<1$ hoặc $r<0$ thì báo lỗi, dừng thuật toán.\newline Nếu $m>n+1$ thì đặt $m:=n+1$.
\item[Bước 1:] Trích $m$ mốc gần $\bar x$ nhất theo thuật toán Câu 11a.\newline Đặt $n:=m-1$.\newline Nếu Gauss I thì đặt $q:=\lfloor n/2\rfloor$; nếu Gauss II thì đặt $q:=\lceil n/2\rceil$.\newline Lấy mốc thứ $q$ (đếm từ $0$) làm $x_0$, đánh số lại $i=-q,\dots,n-q$.
\item[Bước 2:] Khởi tạo ma trận $D$, hàng $i=-q,\dots,n-q$, cột $k=0,\dots,n$ ($D_{i,k}=\Delta^ky_i$).\newline Lặp $i=-q,\dots,n-q$: đặt $D_{i,0}:=y_i$.\newline Lặp lồng $k=1,\dots,n$; $i=-q,\dots,n-q-k$: tính $D_{i,k}:=D_{i+1,k-1}-D_{i,k-1}$.
\item[Bước 3:] Lặp $k=0,\dots,n$:\thut nếu Gauss I thì tính $c_k:=D_{-\lfloor k/2\rfloor,k}/k!$ và đặt $z_k:=(-1)^{k+1}\lceil k/2\rceil$ ($z=0,1,-1,2,\dots$);\thut nếu Gauss II thì tính $c_k:=D_{-\lceil k/2\rceil,k}/k!$ và đặt $z_k:=(-1)^{k}\lceil k/2\rceil$ ($z=0,-1,1,-2,\dots$).
\item[Bước 4:] Khởi tạo vector $e$, $b$, mỗi vector $n+1$ phần tử (chỉ số $=$ số mũ của $t$):\thut $(e_0,e_1,\dots,e_n):=(1,0,\dots,0)$;\thut $(b_0,b_1,\dots,b_n):=(c_0,0,\dots,0)$.\newline Lặp $k=1,\dots,n$:
  \begin{con}
  \item[4.1:] Nhân vector $e$ với $(t-z_{k-1})$ theo thuật toán Câu 5a ($e$ thành hệ số của $\omega_k(t)$).
  \item[4.2:] Lặp $l=0,\dots,k$: tính $b_l:=b_l+c_k\,e_l$.
  \end{con}
\item[Bước 5:] Tính $\bar t:=(\bar x-x_0)/h$.\newline Gọi thuật toán Câu 4c với vector $b$, $c:=\bar t$, cấp $k:=r$, được vector $d=(d_0,\dots,d_r)$.\newline Lặp $j=0,\dots,r$: tính $f_j:=d_j/h^j$.\newline Dừng thuật toán, trả về $b$ và $f$.
\end{thuattoan}

\subsection{b) Stirling}
\begin{thuattoan}{Stirling}
\item[Input:] vector $x$, $y$ (mốc cách đều); điểm $\bar x$; cấp đạo hàm $r$; số mốc dùng $m$.
\item[Output:] vector $b=(b_n,\dots,b_0)$ với $P(t)=\sum_lb_lt^l$, $t=\frac{x-x_0}{h}$; vector $f=(f_0,\dots,f_r)$, $f_j\approx f^{(j)}(\bar x)$.
\item[Bước 0:] Nếu $x$, $y$ khác số phần tử thì báo lỗi, dừng thuật toán.\newline Sắp các mốc tăng theo $x$.\newline Nếu hai mốc trùng nhau thì: cùng $y$ giữ một mốc, khác $y$ báo lỗi, dừng thuật toán.\newline Đặt $n:=(\text{số mốc})-1$.\newline Đặt $h:=\dfrac{x_n-x_0}{n}$.\newline Lặp $i=0,\dots,n-1$:\thut nếu $|x_{i+1}-x_i-h|>10^{-6}h$ thì báo ``mốc không cách đều'', dừng thuật toán.\newline Nếu $m$, $r$ không nguyên, $m<1$ hoặc $r<0$ thì báo lỗi, dừng thuật toán.\newline Nếu $m>n+1$ thì đặt $m:=n+1$.\newline Nếu $m$ chẵn thì đặt $m:=m-1$.
\item[Bước 1:] Trích $m$ mốc gần $\bar x$ nhất theo thuật toán Câu 11a.\newline Đặt $n:=m-1$, $s:=n/2$.\newline Lấy mốc chính giữa làm $x_0$, đánh số lại $i=-s,\dots,s$.
\item[Bước 2:] Khởi tạo ma trận $D$, hàng $i=-s,\dots,s$, cột $k=0,\dots,n$ ($D_{i,k}=\Delta^ky_i$).\newline Lặp $i=-s,\dots,s$: đặt $D_{i,0}:=y_i$.\newline Lặp lồng $k=1,\dots,n$; $i=-s,\dots,s-k$: tính $D_{i,k}:=D_{i+1,k-1}-D_{i,k-1}$.
\item[Bước 3:] Lặp $j=1,\dots,s$:\thut tính $c_{2j-1}:=\dfrac{D_{-j,2j-1}+D_{-j+1,2j-1}}{2\,(2j-1)!}$;\thut tính $c_{2j}:=\dfrac{D_{-j,2j}}{(2j)!}$.
\item[Bước 4:] Khởi tạo vector $e$, $\alpha$, $\beta$, mỗi vector $s$ phần tử (chỉ số $=$ số mũ của $u=t^2$; $\alpha$, $\beta$ là hệ số của $P_{\text{lẻ}}(u)$, $P_{\text{chẵn}}(u)$):\thut $(e_0,e_1,\dots,e_{s-1}):=(1,0,\dots,0)$;\thut $(\alpha_0,\dots,\alpha_{s-1}):=(c_1,0,\dots,0)$, $(\beta_0,\dots,\beta_{s-1}):=(c_2,0,\dots,0)$.\newline Lặp $j=1,\dots,s-1$:
  \begin{con}
  \item[4.1:] Nhân vector $e$ với $(u-j^2)$ theo thuật toán Câu 5a ($e$ thành hệ số của $\Omega_j(u)$).
  \item[4.2:] Lặp $l=0,\dots,j$: tính $\alpha_l:=\alpha_l+c_{2j+1}e_l$, $\beta_l:=\beta_l+c_{2j+2}e_l$.
  \end{con}
\item[Bước 5:] Ghép $P(t)=y_0+t\,P_{\text{lẻ}}(t^2)+t^2P_{\text{chẵn}}(t^2)$:\thut đặt $b_0:=y_0$;\thut lặp $l=0,\dots,s-1$: đặt $b_{2l+1}:=\alpha_l$, $b_{2l+2}:=\beta_l$.
\item[Bước 6:] Tính $\bar t:=(\bar x-x_0)/h$.\newline Gọi thuật toán Câu 4c với vector $b$, $c:=\bar t$, cấp $k:=r$, được vector $d=(d_0,\dots,d_r)$.\newline Lặp $j=0,\dots,r$: tính $f_j:=d_j/h^j$.\newline Dừng thuật toán, trả về $b$ và $f$.
\end{thuattoan}

\subsection{c) Bessel}
\begin{thuattoan}{Bessel}
\item[Input:] vector $x$, $y$ (mốc cách đều); điểm $\bar x$; cấp đạo hàm $r$; số mốc dùng $m$.
\item[Output:] vector $b=(b_n,\dots,b_0)$ với $P(v)=\sum_lb_lv^l$, $v=t-\frac12$; vector $f=(f_0,\dots,f_r)$, $f_j\approx f^{(j)}(\bar x)$.
\item[Bước 0:] Nếu $x$, $y$ khác số phần tử thì báo lỗi, dừng thuật toán.\newline Sắp các mốc tăng theo $x$.\newline Nếu hai mốc trùng nhau thì: cùng $y$ giữ một mốc, khác $y$ báo lỗi, dừng thuật toán.\newline Đặt $n:=(\text{số mốc})-1$.\newline Đặt $h:=\dfrac{x_n-x_0}{n}$.\newline Lặp $i=0,\dots,n-1$:\thut nếu $|x_{i+1}-x_i-h|>10^{-6}h$ thì báo ``mốc không cách đều'', dừng thuật toán.\newline Nếu $m$, $r$ không nguyên, $m<1$ hoặc $r<0$ thì báo lỗi, dừng thuật toán.\newline Nếu $m>n+1$ thì đặt $m:=n+1$.\newline Nếu $m$ lẻ thì đặt $m:=m-1$.
\item[Bước 1:] Trích $m$ mốc gần $\bar x$ nhất theo thuật toán Câu 11a.\newline Đặt $n:=m-1$, $s:=\frac{n+1}{2}$.\newline Lấy mốc thứ $s-1$ (đếm từ $0$) làm $x_0$ ($x_0$, $x_1$ là hai mốc chính giữa), đánh số lại $i=-(s-1),\dots,s$.
\item[Bước 2:] Khởi tạo ma trận $D$, hàng $i=-(s-1),\dots,s$, cột $k=0,\dots,n$ ($D_{i,k}=\Delta^ky_i$).\newline Lặp $i=-(s-1),\dots,s$: đặt $D_{i,0}:=y_i$.\newline Lặp lồng $k=1,\dots,n$; $i=-(s-1),\dots,s-k$: tính $D_{i,k}:=D_{i+1,k-1}-D_{i,k-1}$.
\item[Bước 3:] Lặp $j=0,\dots,s-1$:\thut tính $c_{2j}:=\dfrac{D_{-j,2j}+D_{-j+1,2j}}{2\,(2j)!}$;\thut tính $c_{2j+1}:=\dfrac{D_{-j,2j+1}}{(2j+1)!}$.
\item[Bước 4:] Khởi tạo vector $e$, $\alpha$, $\beta$, mỗi vector $s$ phần tử (chỉ số $=$ số mũ của $u=v^2$; $\alpha$, $\beta$ là hệ số của $P_{\text{lẻ}}(u)$, $P_{\text{chẵn}}(u)$):\thut $(e_0,e_1,\dots,e_{s-1}):=(1,0,\dots,0)$;\thut $(\alpha_0,\dots,\alpha_{s-1}):=(c_1,0,\dots,0)$, $(\beta_0,\dots,\beta_{s-1}):=(c_0,0,\dots,0)$.\newline Lặp $j=1,\dots,s-1$:
  \begin{con}
  \item[4.1:] Nhân vector $e$ với $\big(u-(j-\frac12)^2\big)$ theo thuật toán Câu 5a ($e$ thành hệ số của $\Omega_j(u)$).
  \item[4.2:] Lặp $l=0,\dots,j$: tính $\alpha_l:=\alpha_l+c_{2j+1}e_l$, $\beta_l:=\beta_l+c_{2j}e_l$.
  \end{con}
\item[Bước 5:] Ghép $P(v)=P_{\text{chẵn}}(v^2)+v\,P_{\text{lẻ}}(v^2)$:\thut lặp $l=0,\dots,s-1$: đặt $b_{2l}:=\beta_l$, $b_{2l+1}:=\alpha_l$.
\item[Bước 6:] Tính $\bar v:=(\bar x-x_0)/h-\frac12$.\newline Gọi thuật toán Câu 4c với vector $b$, $c:=\bar v$, cấp $k:=r$, được vector $d=(d_0,\dots,d_r)$.\newline Lặp $j=0,\dots,r$: tính $f_j:=d_j/h^j$.\newline Dừng thuật toán, trả về $b$ và $f$.
\end{thuattoan}

%=====================================================================
\section{Câu 16. Trích xuất 9 điểm và nội suy Stirling tại $\bar x=7.83$}

\subsection{a) Trích xuất}
Dữ liệu \texttt{data Câu 11.xlsx} đã xử lý ở câu 11b (141 mốc, $h=0.1$). Thuật toán câu 11a ($\bar x=7.83$, $k=9$): <<c16_ket_trich>> Số mốc trái bằng số mốc phải nên mốc tâm nằm chính giữa: $x_0=<<c16_x0>>$, đánh số lại $i=<<c16_p>>..<<c16_pn>>$:
<<c16_bang_xy>>

\subsection{b) Đa thức Stirling, $f(\bar x)$, $f'(\bar x)$}
\pp{Phương pháp.} 9 mốc cách đều (số lẻ) nên dùng Stirling, mốc trung tâm $x_0=<<c16_x0>>$ ($i=<<c16_p>>..<<c16_pn>>$), $h=0.1$, $t=\frac{x-x_0}{h}$:
\[ P(x_0+th)=y_0+t\,\frac{\Delta y_{-1}+\Delta y_0}{2}+\frac{t^2}{2!}\Delta^2y_{-1}+\frac{t(t^2-1)}{3!}\cdot\frac{\Delta^3y_{-2}+\Delta^3y_{-1}}{2}+\frac{t^2(t^2-1)}{4!}\Delta^4y_{-2}+\cdots \]

\pp{Bảng sai phân} ($\Delta^0y_i=y_i$, $\Delta^ky_i=\Delta^{k-1}y_{i+1}-\Delta^{k-1}y_i$):
<<c16_bang_sp>>
\pp{Hệ số:}
<<c16_C>>
\pp{Bảng tích theo $u=t^2$} ($\Omega_j(u)=\Omega_{j-1}(u)\cdot(u-j^2)$: $A_{0,0}=1$, $A_{j,l}=A_{j-1,l-1}-j^2A_{j-1,l}$):
<<c16_bang_tich>>
<<c16_CA>>
Ghép hệ số theo $P(t)=y_0+t\,P_{\text{lẻ}}(t^2)+t^2P_{\text{chẵn}}(t^2)$:
<<c16_Pt>>
với $t=(x-7.8)/0.1$.

\pp{Tại $\bar t=\frac{\bar x-x_0}{h}=<<c16_t>>$:} Horner chia lặp $P_8$ tại $\bar t$ (Câu 4c):
\[ f(7.83)\approx P_8(\bar t)=<<c16_f>>,\qquad f'(7.83)\approx\frac{P_8'(\bar t)}{h}=\frac{<<c16_dP>>}{0.1}=<<c16_df>>. \]

%=====================================================================
\section{Câu 17. Trích xuất 8 điểm và nội suy Bessel tại $\bar x=7.83$}

\subsection{a) Trích xuất}
Thuật toán câu 11a ($\bar x=7.83$, $k=8$, dữ liệu câu 11b): <<c17_ket_trich>> $\bar x$ nằm giữa hai mốc chính giữa $x_0=<<c17_x0>>$ (mốc tâm) và $x_1=7.9$, đánh số lại $i=<<c17_p>>..<<c17_pn>>$:
<<c17_bang_xy>>

\subsection{b) Đa thức Bessel, $f(\bar x)$, $f'(\bar x)$}
\pp{Phương pháp.} 8 mốc cách đều (số chẵn) nên dùng Bessel, hai mốc chính giữa $x_0=<<c17_x0>>$, $x_1=7.9$ ($i=<<c17_p>>..<<c17_pn>>$), $h=0.1$, $t=\frac{x-x_0}{h}$:
\begin{align*}
P(x_0+th)&=\frac{y_0+y_1}{2}+\Big(t-\frac12\Big)\Delta y_0+\frac{t(t-1)}{2!}\cdot\frac{\Delta^2y_{-1}+\Delta^2y_0}{2}+\frac{(t-\frac12)t(t-1)}{3!}\Delta^3y_{-1}\\
&\quad+\frac{(t+1)t(t-1)(t-2)}{4!}\cdot\frac{\Delta^4y_{-2}+\Delta^4y_{-1}}{2}+\cdots
\end{align*}

\pp{Bảng sai phân} ($\Delta^0y_i=y_i$, $\Delta^ky_i=\Delta^{k-1}y_{i+1}-\Delta^{k-1}y_i$):
<<c17_bang_sp>>
\pp{Hệ số:}
<<c17_C>>
\pp{Bảng tích theo $u=v^2$} ($\Omega_j(u)=\Omega_{j-1}(u)\cdot\big(u-(j-\frac12)^2\big)$: $A_{0,0}=1$, $A_{j,l}=A_{j-1,l-1}-(j-\tfrac12)^2A_{j-1,l}$):
<<c17_bang_tich>>
<<c17_CA>>
Ghép hệ số theo $P(v)=P_{\text{chẵn}}(v^2)+v\,P_{\text{lẻ}}(v^2)$:
<<c17_Pt>>
với $v=t-\frac12=(x-<<c17_s0>>)/0.1$.

\pp{Tại $\bar t=\frac{\bar x-x_0}{h}=<<c17_t>>$, $\bar v=\bar t-\frac12=<<c17_v>>$:} Horner chia lặp $P_7$ tại $\bar v$ (Câu 4c):
\[ f(7.83)\approx P_7(\bar v)=<<c17_f>>,\qquad f'(7.83)\approx\frac{P_7'(\bar v)}{h}=\frac{<<c17_dP>>}{0.1}=<<c17_df>>. \]

%=====================================================================
\section{Câu 18. Thuật toán cho bài toán nội suy ngược}

\subsection{a) Tất cả các khoảng cách ly của $f(x)=\bar y$}
\begin{thuattoan}{Khoảng cách ly}
\item[Input:] vector $x$, $y$ ($n+1$ mốc); giá trị $\bar y$.
\item[Output:] danh sách $N$ các mốc là nghiệm; danh sách $K$ các khoảng cách ly $(x_i,x_{i+1})$.
\item[Bước 0:] Nếu $x$, $y$ khác số phần tử thì báo lỗi, dừng thuật toán.\newline Sắp các mốc tăng theo $x$.\newline Nếu hai mốc trùng nhau thì: cùng $y$ giữ một mốc, khác $y$ báo lỗi, dừng thuật toán.\newline Đặt $n:=(\text{số mốc})-1$.
\item[Bước 1:] Đặt $N:=\varnothing$, $K:=\varnothing$.\newline Lặp $i=0,\dots,n$: tính $g_i:=y_i-\bar y$.
\item[Bước 2:] Lặp $i=0,\dots,n$:\thut nếu $g_i=0$ thì thêm $x_i$ vào $N$;\thut nếu $i<n$ và $g_i\,g_{i+1}<0$ thì thêm $(x_i,x_{i+1})$ vào $K$.
\item[Bước 3:] Nếu $N$ và $K$ đều rỗng thì báo ``phương trình không có nghiệm trong $[x_0,x_n]$''.\newline Dừng thuật toán, trả về $N$ và $K$.
\end{thuattoan}

\subsection{b) Các khoảng đơn điệu}
\begin{thuattoan}{Khoảng đơn điệu}
\item[Input:] vector $x$, $y$ ($n+1$ mốc).
\item[Output:] danh sách các đoạn $[x_a,x_b]$ kèm chiều tăng / giảm / hằng.
\item[Bước 0:] Nếu $x$, $y$ khác số phần tử thì báo lỗi, dừng thuật toán.\newline Sắp các mốc tăng theo $x$.\newline Nếu hai mốc trùng nhau thì: cùng $y$ giữ một mốc, khác $y$ báo lỗi, dừng thuật toán.\newline Đặt $n:=(\text{số mốc})-1$.\newline Nếu $n<1$ thì báo lỗi (cần ít nhất 2 mốc), dừng thuật toán.
\item[Bước 1:] Lặp $i=0,\dots,n-1$: tính $s_i:=\operatorname{sign}(y_{i+1}-y_i)$.\newline Đặt $a:=0$.
\item[Bước 2:] Lặp $i=1,\dots,n-1$:\thut nếu $s_i\ne s_{i-1}$ thì ghi đoạn $[x_a,x_i]$ với chiều $s_{i-1}$ và đặt $a:=i$.
\item[Bước 3:] Ghi đoạn $[x_a,x_n]$ với chiều $s_{n-1}$.\newline Dừng thuật toán, trả về danh sách các đoạn.
\end{thuattoan}

\subsection{c) Phương pháp hàm ngược}
\begin{thuattoan}{Hàm ngược}
\item[Input:] vector $x$, $y$; $\bar y$; khoảng cách ly $(x_k,x_{k+1})$ nằm trong đoạn đơn điệu $[x_a,x_b]$ (câu a, b); số mốc $m$.
\item[Output:] nghiệm gần đúng $\bar x$.
\item[Bước 0:] Nếu $x$, $y$ khác số phần tử thì báo lỗi, dừng thuật toán.\newline Sắp các mốc tăng theo $x$.\newline Nếu hai mốc trùng nhau thì: cùng $y$ giữ một mốc, khác $y$ báo lỗi, dừng thuật toán.\newline Đặt $n:=(\text{số mốc})-1$.\newline Nếu $m$ không nguyên hoặc $m<2$ thì báo lỗi, dừng thuật toán.
\item[Bước 1:] Đặt $L:=k$, $R:=k+1$ (các mốc đã chọn là $x_L,\dots,x_R$).
\item[Bước 2:] Lặp $s=3,\dots,m$ (mỗi lần thêm mốc kề trong $[x_a,x_b]$ có $|y-\bar y|$ nhỏ hơn):\thut nếu $L=a$ và $R=b$ thì dừng lặp (đoạn đơn điệu hết mốc);\thut ngược lại nếu $L=a$ thì $R:=R+1$;\thut ngược lại nếu $R=b$ thì $L:=L-1$;\thut ngược lại nếu $|y_{L-1}-\bar y|\le|y_{R+1}-\bar y|$ thì $L:=L-1$, còn không thì $R:=R+1$.
\item[Bước 3:] Đặt $u:=(y_L,\dots,y_R)$, $v:=(x_L,\dots,x_R)$ (đánh số lại $0,\dots,R-L$).\newline Nếu đoạn giảm thì đảo thứ tự $u$, $v$ để $u_0<u_1<\cdots$.
\item[Bước 4:] Khởi tạo ma trận $T$ cỡ $(R-L+1)\times(R-L+1)$ (tỷ sai phân của $v$ theo $u$).\newline Lặp $i=0,\dots,R-L$: đặt $T_{i,0}:=v_i$.\newline Lặp lồng $k=1,\dots,R-L$; $i=k,\dots,R-L$: tính $T_{i,k}:=\dfrac{T_{i,k-1}-T_{i-1,k-1}}{u_i-u_{i-k}}$.\newline Lặp $j=0,\dots,R-L$: đặt $c_j:=T_{j,j}$.
\item[Bước 5:] Đặt $\omega:=1$, $S:=c_0$.\newline Lặp $j=1,\dots,R-L$:\thut tính $\omega:=\omega\,(\bar y-u_{j-1})$;\thut tính $S:=S+c_j\,\omega$.
\item[Bước 6:] Đặt $\bar x:=S$.\newline Dừng thuật toán, trả về $\bar x$ (số hạng cuối $|c_{R-L}\,\omega|$ là ước lượng sai số).
\end{thuattoan}

\subsection{d) Phương pháp lặp}
\begin{thuattoan}{Lặp Newton tiến / lùi}
\item[Input:] vector $x$, $y$ ($n+1$ mốc cách đều bước $h$); $\bar y$; khoảng cách ly $(x_k,x_{k+1})$; số mốc $m$; sai số $\varepsilon$; số lần lặp tối đa $N$.
\item[Output:] nghiệm gần đúng $\bar x$, hoặc thông báo ``Không hội tụ''.
\item[Bước 0:] Nếu $x$, $y$ khác số phần tử thì báo lỗi, dừng thuật toán.\newline Sắp các mốc tăng theo $x$.\newline Nếu hai mốc trùng nhau thì: cùng $y$ giữ một mốc, khác $y$ báo lỗi, dừng thuật toán.\newline Đặt $n:=(\text{số mốc})-1$.\newline Đặt $h:=\dfrac{x_n-x_0}{n}$.\newline Lặp $i=0,\dots,n-1$:\thut nếu $|x_{i+1}-x_i-h|>10^{-6}h$ thì dùng phương pháp hàm ngược (câu c), dừng thuật toán.\newline Nếu $m$ không nguyên hoặc $m<2$ thì báo lỗi, dừng thuật toán.\newline Nếu $k+m-1>n$ và $k+2-m<0$ thì đặt $m:=\max(n-k+1,\,k+2)$.
\item[Bước 1:] Chọn chiều:\thut nếu $k+m-1\le n$ thì lặp tiến, đặt $s:=k$;\thut trái lại lặp lùi, đặt $s:=k+1$.
\item[Bước 2:] Lập bảng sai phân $T$ ($T_{i,0}:=y_i$, $T_{i,k}:=T_{i,k-1}-T_{i-1,k-1}$), đánh số lại $0..m-1$:\thut nếu tiến thì của các mốc $x_k,\dots,x_{k+m-1}$, đặt $d_j:=T_{j,j}=\Delta^jy_k$;\thut nếu lùi thì của các mốc $x_{k+2-m},\dots,x_{k+1}$, đặt $d_j:=T_{m-1,j}=\nabla^jy_{k+1}$ ($j=0..m-1$).
\item[Bước 3:] Nếu tiến thì $\varphi(t):=\dfrac{\bar y-d_0}{d_1}-\sum_{j=2}^{m-1}\dfrac{d_j}{j!\,d_1}\,t(t-1)\cdots(t-j+1)$.\newline Nếu lùi thì $\varphi(t):=\dfrac{\bar y-d_0}{d_1}-\sum_{j=2}^{m-1}\dfrac{d_j}{j!\,d_1}\,t(t+1)\cdots(t+j-1)$.
\item[Bước 4:] Đặt $t_0:=\dfrac{\bar y-d_0}{d_1}$, $l:=0$.
\item[Bước 5:] Tính $t_{l+1}:=\varphi(t_l)$.
\item[Bước 6:] Nếu $h\,|t_{l+1}-t_l|<\varepsilon$ thì đặt $\bar x:=x_s+t_{l+1}h$, dừng thuật toán.
\item[Bước 7:] Nếu $l+1=N$ thì báo ``Không hội tụ'', dừng thuật toán.\newline Trái lại đặt $l:=l+1$, quay lại bước 5.
\end{thuattoan}

%=====================================================================
\section{Câu 19. Giải phương trình $f(x)=\bar y$ từ file dữ liệu}

\pp{Dữ liệu.} \texttt{data Câu 19.xlsx}: $<<c19_N>>$ mốc $x=1;1.1;\dots;15$ cách đều $h=0.1$ (dòng đầu file trống, bỏ qua). Các khoảng đơn điệu (câu 18b):
<<c19_dd>>

\subsection{a) Phương pháp hàm ngược, $\bar y=2.5$}
\pp{Khoảng cách ly} (câu 18a):
<<c19a_cl>>
Chỉ có một nghiệm, trong $(x_{<<c19a_k>>};x_{<<c19a_k>>+1})=(<<c19a_xk>>;\ <<c19a_xk1>>)\subset<<c19a_doan>>$ — đoạn $f$ tăng chặt nên có hàm ngược.

\pp{Chọn 6 mốc} có $|y_j-\bar y|$ nhỏ nhất trong đoạn đơn điệu:
<<c19a_chon>>
\pp{Bảng tỷ sai phân của $x$ theo $y$} ($u=y$ tăng dần; đường chéo khung $=c_j$):
<<c19a_tsp>>
\pp{Giá trị tại $\bar y$:} $\bar x=\sum_jc_j\,\omega_j(\bar y)$, $\omega_j=\prod_{l<j}(\bar y-u_l)$:
<<c19a_tong>>
\begin{ketluan}
$f(x)=2.5$ có đúng một nghiệm trong $[1;15]$: $\bar x\approx<<c19a_x>>$ (số hạng cuối $<<c19a_cuoi>>$ — ước lượng sai số).
\end{ketluan}

\subsection{b) Phương pháp lặp, $\bar y=3.15$}
\pp{Khoảng cách ly:}
<<c19b_cl>>
Có $<<c19b_so>>$ nghiệm. Mỗi khoảng dùng $m=6$ mốc, $\varepsilon=10^{-6}$, $N=100$. <<c19b_chung>>
<<c19b_khoi>>
\begin{ketluan}
$f(x)=3.15$ có 3 nghiệm trong $[1;15]$: $\bar x=<<c19b_x>>$.
\end{ketluan}

%=====================================================================
\section{Câu 20. Hàm ghép trơn (spline)}

\subsection{a) Xây dựng công thức}
Spline cấp $p$: $S(x)=S_k(x)$ trên $[x_k,x_{k+1}]$, $S_k$ đa thức bậc $p$, $S\in C^{p-1}[x_0,x_n]$, $S(x_i)=y_i$; $h_k=x_{k+1}-x_k$; cần thêm $p-1$ điều kiện biên.

\pp{Spline cấp 3.} Đặt $\sigma_k=S''(x_k)$. $S_k''$ bậc nhất nên $S_k''(x)=\dfrac{\sigma_k(x_{k+1}-x)+\sigma_{k+1}(x-x_k)}{h_k}$; tích phân hai lần:
\[ S_k(x)=\frac{\sigma_k(x_{k+1}-x)^3+\sigma_{k+1}(x-x_k)^3}{6h_k}+\alpha_k(x-x_k)+\beta_k(x_{k+1}-x). \]
$S_k(x_k)=y_k\Rightarrow\beta_k=\dfrac{y_k}{h_k}-\dfrac{\sigma_kh_k}{6}$; \ $S_k(x_{k+1})=y_{k+1}\Rightarrow\alpha_k=\dfrac{y_{k+1}}{h_k}-\dfrac{\sigma_{k+1}h_k}{6}$. Đặt $s_k=\dfrac{y_{k+1}-y_k}{h_k}$:
\[ S_k'(x)=\frac{-\sigma_k(x_{k+1}-x)^2+\sigma_{k+1}(x-x_k)^2}{2h_k}+s_k-\frac{(\sigma_{k+1}-\sigma_k)h_k}{6}, \]
\[ S_k'(x_{k+1})=s_k+\frac{h_k}{6}\sigma_k+\frac{h_k}{3}\sigma_{k+1},\qquad S_{k+1}'(x_{k+1})=s_{k+1}-\frac{h_{k+1}}{3}\sigma_{k+1}-\frac{h_{k+1}}{6}\sigma_{k+2}. \]
Trơn cấp 1 tại $x_{k+1}$ cho $n-1$ phương trình ($k=0..n-2$):
\[ \boxed{\frac{h_k}{6}\sigma_k+\frac{h_k+h_{k+1}}{3}\sigma_{k+1}+\frac{h_{k+1}}{6}\sigma_{k+2}=s_{k+1}-s_k} \]
cộng 2 điều kiện biên: $S''(x_0)=A\Rightarrow\sigma_0=A$; $S'(x_0)=A\Rightarrow\frac{h_0}{3}\sigma_0+\frac{h_0}{6}\sigma_1=s_0-A$; $S''(x_n)=B\Rightarrow\sigma_n=B$; $S'(x_n)=B\Rightarrow\frac{h_{n-1}}{6}\sigma_{n-1}+\frac{h_{n-1}}{3}\sigma_n=B-s_{n-1}$. Hệ ba đường chéo, chéo trội. Khai triển được dạng chính tắc $S_k=a_kx^3+b_kx^2+c_kx+d_k$:
\[ a_k=\frac{\sigma_{k+1}-\sigma_k}{6h_k},\ b_k=\frac{\sigma_kx_{k+1}-\sigma_{k+1}x_k}{2h_k},\ c_k=\frac{-\sigma_kx_{k+1}^2+\sigma_{k+1}x_k^2}{2h_k}+s_k-\frac{(\sigma_{k+1}-\sigma_k)h_k}{6}, \]
\[ d_k=\frac{\sigma_kx_{k+1}^3-\sigma_{k+1}x_k^3}{6h_k}+\frac{y_kx_{k+1}-y_{k+1}x_k}{h_k}+\frac{(\sigma_{k+1}x_k-\sigma_kx_{k+1})h_k}{6}. \]

\pp{Spline cấp 4.} Đặt $m_k=S'(x_k)$, $\sigma_k=S''(x_k)$, $M_k=S'''(x_k)$ ($S\in C^3$ nên các giá trị này chung cho hai mảnh kề). $S_k'''$ bậc nhất: $S_k'''(x)=M_k+\frac{M_{k+1}-M_k}{h_k}u$, $u=x-x_k$. Tích phân ba lần với giá trị ban đầu tại $x_k$ (khai triển Taylor):
\[ S_k(x)=y_k+m_ku+\frac{\sigma_k}{2}u^2+\frac{M_k}{6}u^3+\frac{M_{k+1}-M_k}{24h_k}u^4,\qquad u=x-x_k. \]
Liên tục của $S,S',S''$ tại $x_{k+1}$ ($S'''$ liên tục sẵn), $k=0..n-1$:
\[ \boxed{\begin{aligned}
&y_k+m_kh_k+\frac{\sigma_k}{2}h_k^2+\frac{3M_k+M_{k+1}}{24}h_k^3=y_{k+1},\\
&m_k+\sigma_kh_k+\frac{2M_k+M_{k+1}}{6}h_k^2=m_{k+1},\\
&\sigma_k+\frac{M_k+M_{k+1}}{2}h_k=\sigma_{k+1}.
\end{aligned}} \]
$3n$ phương trình, $3(n+1)$ ẩn $(m_k,\sigma_k,M_k)$ $\Rightarrow$ cần 3 điều kiện biên (vd $S'(x_0),S''(x_0),S'(x_n)$). Dạng chính tắc: khai triển $u=x-x_k$ bằng Horner nhân.

\subsection{b) Thuật toán spline cấp $p\le4$}
\begin{thuattoan}{Spline cấp 1}
\item[Input:] $(x_i,y_i)$, $i=0..n$.
\item[Output:] $S_k(x)=a_kx+b_k$, $k=0..n-1$.
\item[Bước 0:] Sắp các mốc tăng theo $x$.\newline Nếu hai mốc trùng nhau thì báo lỗi, dừng thuật toán.
\item[Bước 1:] Lặp $k=0,\dots,n-1$:\thut tính $h_k:=x_{k+1}-x_k$;\thut tính $a_k:=\dfrac{y_{k+1}-y_k}{h_k}$, $b_k:=\dfrac{y_kx_{k+1}-y_{k+1}x_k}{h_k}$.\newline Dừng thuật toán, trả về các $S_k$.
\end{thuattoan}

\begin{thuattoan}{Spline cấp 2}
\item[Input:] $(x_i,y_i)$, $i=0..n$; một điều kiện biên $S'(x_0)=A$ hoặc $S'(x_n)=B$.
\item[Output:] $S_k(x)=a_kx^2+b_kx+c_k$, $k=0..n-1$.
\item[Bước 0:] Sắp các mốc tăng theo $x$.\newline Nếu hai mốc trùng nhau thì báo lỗi, dừng thuật toán.\newline Nếu số điều kiện biên khác $1$ hoặc có điều kiện không đặt tại $x_0$, $x_n$ thì báo lỗi, dừng thuật toán.
\item[Bước 1:] Lặp $k=0,\dots,n-1$: tính $h_k:=x_{k+1}-x_k$ và $\gamma_k:=\dfrac{2(y_{k+1}-y_k)}{h_k}$.
\item[Bước 2:] Tính $m_k=S'(x_k)$:
  \begin{con}
  \item[TH1:] Nếu cho $S'(x_0)=A$ thì:\thut đặt $m_0:=A$;\thut lặp $k=0,\dots,n-1$: tính $m_{k+1}:=\gamma_k-m_k$.
  \item[TH2:] Nếu cho $S'(x_n)=B$ thì:\thut đặt $m_n:=B$;\thut lặp $k=n-1,\dots,0$: tính $m_k:=\gamma_k-m_{k+1}$.
  \end{con}
\item[Bước 3:] Lặp $k=0,\dots,n-1$: tính
  \[ a_k:=\frac{m_{k+1}-m_k}{2h_k},\ \ b_k:=\frac{m_kx_{k+1}-m_{k+1}x_k}{h_k},\ \ c_k:=\frac{-m_kx_{k+1}^2+m_{k+1}x_k^2}{2h_k}+y_k+\frac{m_kh_k}{2}. \]
  Dừng thuật toán, trả về các $S_k$.
\end{thuattoan}

\begin{thuattoan}{Spline cấp 3}
\item[Input:] $(x_i,y_i)$, $i=0..n$; hai điều kiện biên dạng $S'(\xi)=v$ hoặc $S''(\xi)=v$, $\xi\in\{x_0,x_n\}$.
\item[Output:] $S_k(x)=a_kx^3+b_kx^2+c_kx+d_k$, $k=0..n-1$.
\item[Bước 0:] Sắp các mốc tăng theo $x$.\newline Nếu hai mốc trùng nhau thì báo lỗi, dừng thuật toán.\newline Nếu số điều kiện biên khác $2$ hoặc có điều kiện không đặt tại $x_0$, $x_n$ thì báo lỗi, dừng thuật toán.
\item[Bước 1:] Lặp $k=0,\dots,n-1$: tính $h_k:=x_{k+1}-x_k$ và $s_k:=\dfrac{y_{k+1}-y_k}{h_k}$.
\item[Bước 2:] Lập hệ $n+1$ phương trình ẩn $\sigma_0,\dots,\sigma_n$:\thut $n-1$ phương trình trơn (khung, câu a);\thut 2 phương trình biên (câu a).
\item[Bước 3:] Nếu mỗi đầu có một điều kiện thì hệ có dạng ba đường chéo $\ell_i\sigma_{i-1}+d_i\sigma_i+u_i\sigma_{i+1}=r_i$, giải bằng truy đuổi:
  \begin{con}
  \item[3.1:] Đặt $P_0:=-u_0/d_0$, $Q_0:=r_0/d_0$.\newline Lặp $i=1,\dots,n$:\thut tính $\mu:=d_i+\ell_iP_{i-1}$; nếu $\mu=0$ thì báo hệ suy biến, dừng thuật toán;\thut tính $P_i:=-u_i/\mu$, $Q_i:=(r_i-\ell_iQ_{i-1})/\mu$.
  \item[3.2:] Đặt $\sigma_n:=Q_n$.\newline Lặp $i=n-1,\dots,0$: tính $\sigma_i:=P_i\sigma_{i+1}+Q_i$.\newline Sang bước 5.
  \end{con}
\item[Bước 4:] Nếu hai điều kiện cùng một đầu thì giải hệ bằng khử Gauss chọn trụ.
\item[Bước 5:] Lặp $k=0,\dots,n-1$: tính $a_k,b_k,c_k,d_k$ theo câu a.\newline Dừng thuật toán, trả về các $S_k$.
\end{thuattoan}

\begin{thuattoan}{Spline cấp 4}
\item[Input:] $(x_i,y_i)$, $i=0..n$; ba điều kiện biên $S^{(r)}(\xi)=v$, $r\in\{1,2,3\}$, $\xi\in\{x_0,x_n\}$.
\item[Output:] hệ số của $S_k(x)$ (bậc 4, từ bậc cao xuống), $k=0..n-1$.
\item[Bước 0:] Sắp các mốc tăng theo $x$.\newline Nếu hai mốc trùng nhau thì báo lỗi, dừng thuật toán.\newline Nếu số điều kiện biên khác $3$ hoặc có điều kiện không đặt tại $x_0$, $x_n$ thì báo lỗi, dừng thuật toán.
\item[Bước 1:] Lặp $k=0,\dots,n-1$: tính $h_k:=x_{k+1}-x_k$.
\item[Bước 2:] Lập hệ $3(n+1)$ phương trình ẩn $m_k,\sigma_k,M_k$ ($k=0..n$):\thut $3n$ phương trình liên tục (khung, câu a);\thut 3 phương trình biên: $S'(\xi)=v\Rightarrow m=v$, $S''(\xi)=v\Rightarrow\sigma=v$, $S'''(\xi)=v\Rightarrow M=v$ tại mốc $\xi$.
\item[Bước 3:] Giải hệ bằng khử Gauss chọn trụ.\newline Nếu trụ bằng $0$ thì báo ``điều kiện biên không xác định duy nhất spline'', dừng thuật toán.
\item[Bước 4:] Lặp $k=0,\dots,n-1$:
  \begin{con}
  \item[4.1:] Đặt $e:=\Big(\dfrac{M_{k+1}-M_k}{24h_k},\ \dfrac{M_k}{6},\ \dfrac{\sigma_k}{2},\ m_k,\ y_k\Big)$ — hệ số của $S_k$ theo $u=x-x_k$.
  \item[4.2:] Đặt $S_k:=(e_4)$.\newline Lặp $j=3,\dots,0$: nhân $S_k$ với $(x-x_k)$ theo thuật toán Câu 5a rồi cộng $e_j$ vào hệ số tự do.
  \end{con}
\item[Bước 5:] Dừng thuật toán, trả về các $S_k$.
\end{thuattoan}

%=====================================================================
\section{Câu 21. Spline cho bảng dữ liệu}

\pp{Xử lý dữ liệu.} Bảng đề cho $x=0;\,0.5236;\,1.5708;\,1.0472;\,3.1416$ \textbf{không tăng dần} (cặp $(1.0472;-0.866)$ đứng sau $(1.5708;0)$). Giữ nguyên số liệu, sắp lại theo $x$:
<<c21_xy>>
<<c21_h>>

\subsection{a) Spline cấp 2 với $S'(0)=0$}
\pp{Phương pháp.} Thuật toán spline cấp 2 (Câu 20b), TH1: $m_0=S'(0)=0$, truy hồi xuôi $m_{k+1}=\gamma_k-m_k$.
<<c21a_m>>
$m=<<c21a_mv>>$.
<<c21a_he>>
\begin{ketluan}
<<c21a_S>>
\end{ketluan}

\subsection{b) Spline cấp 3 với $S''(0)=-1$, $S''(3.1416)=1$}
\pp{Đề.} Đề ghi ``$S''(0)=-1,\ S''(0)=1$'' (lặp mốc 0); hiểu là $S''(0)=-1$, $S''(3.1416)=1$.

\pp{Phương pháp.} Thuật toán spline cấp 3 (Câu 20b): hệ ba đường chéo $\ell_i\sigma_{i-1}+d_i\sigma_i+u_i\sigma_{i+1}=r_i$ theo $\sigma_k=S''(x_k)$, hàng biên $\sigma_0=-1$, $\sigma_4=1$; giải bằng truy đuổi:
<<c21b_duoi>>
Với $\sigma$ trên, hệ số chính tắc tính theo câu 20a:
\begin{ketluan}
<<c21b_S>>
\end{ketluan}

\subsection{c) Spline cấp 3 với $S'(0)=S'(3.1416)=0$}
\pp{Phương pháp.} Như câu b, hàng biên thay bằng
$\frac{h_0}{3}\sigma_0+\frac{h_0}{6}\sigma_1=s_0-0$, \quad $\frac{h_3}{6}\sigma_3+\frac{h_3}{3}\sigma_4=0-s_3$:
<<c21c_duoi>>
\begin{ketluan}
<<c21c_S>>
\end{ketluan}

%=====================================================================
\section{Câu 22. Thuật toán phương pháp bình phương tối thiểu}

\subsection{a) Hàm thực nghiệm với tập cơ sở $\{\varphi_j\}_{j=1}^m$ và sai số trung bình phương}
\begin{thuattoan}{Bình phương tối thiểu tuyến tính theo tham số}
\item[Input:] vector $x$, $y$ ($n$ điểm $(x_i,y_i)$, $i=1..n$); cơ sở $\varphi_1,\dots,\varphi_m$.
\item[Output:] vector hệ số $a=(a_1,\dots,a_m)$ của $g=\sum_ja_j\varphi_j$; sai số trung bình phương $\sigma$.
\item[Bước 0:] Loại các điểm làm $\varphi_j(x_i)$ không xác định.\newline Nếu còn $n<m$ điểm thì báo lỗi, dừng thuật toán.
\item[Bước 1:] Khởi tạo ma trận $\Phi$ cỡ $n\times m$.\newline Lặp lồng $i=1,\dots,n$; $j=1,\dots,m$: tính $\Phi_{ij}:=\varphi_j(x_i)$.
\item[Bước 2:] Khởi tạo ma trận $M$ cỡ $m\times m$ và vector $b$ có $m$ phần tử.\newline Lặp lồng $j=1,\dots,m$; $k=1,\dots,m$: tính $M_{jk}:=\sum_{i=1}^n\Phi_{ij}\Phi_{ik}$.\newline Lặp $j=1,\dots,m$: tính $b_j:=\sum_{i=1}^n\Phi_{ij}y_i$.
\item[Bước 3:] Giải hệ chuẩn $Ma=b$ bằng phân tích Cholesky $M=LL^{T}$ ($M$ đối xứng, xác định dương; $L$ tam giác dưới): giải $Lu=b$ (thế xuôi) rồi $L^{T}a=u$ (thế ngược).\newline Nếu không phân tích được ($M$ không xác định dương) thì báo ``các hàm cơ sở phụ thuộc tuyến tính trên dữ liệu'', dừng thuật toán.
\item[Bước 4:] Tính $\sigma:=\sqrt{\frac1n\sum_{i=1}^n\big(\sum_ja_j\Phi_{ij}-y_i\big)^2}$.\newline Dừng thuật toán, trả về $a$ và $\sigma$.
\end{thuattoan}

\subsection{b) Dạng $y=a\,e^{b_1\varphi_1(x)+b_2\varphi_2(x)}$ và $y=a\,x^b$}
\begin{thuattoan}{Dạng mũ $y=a\,e^{b_1\varphi_1(x)+b_2\varphi_2(x)}$}
\item[Input:] vector $x$, $y$ ($n$ điểm); cơ sở $\varphi_1,\varphi_2$ (không phải hàm hằng).
\item[Output:] $a$, $b_1$, $b_2$ và sai số trung bình phương $\sigma$.
\item[Bước 0:] Loại các điểm làm $\varphi_1(x_i)$, $\varphi_2(x_i)$ không xác định.\newline Đếm $n_+$ (số $y_i>0$) và $n_-$ (số $y_i<0$); đặt $s:=1$ nếu $n_+\ge n_-$, ngược lại $s:=-1$ ($s$ là dấu của $a$).\newline Lặp $i=1,\dots,n$:\thut nếu $s\,y_i\le0$ ($y_i=0$ hoặc trái dấu với $s$) thì loại điểm $i$, báo lý do.\newline Nếu còn ít hơn 3 điểm thì báo lỗi, dừng thuật toán.
\item[Bước 1:] Lặp $i=1,\dots,n$: tính $z_i:=\ln(s\,y_i)$.
\item[Bước 2:] Tìm $(c_0,c_1,c_2)$ theo câu a với dữ liệu $(x_i,z_i)$ và cơ sở $(1,\varphi_1,\varphi_2)$.
\item[Bước 3:] Đặt $a:=s\,e^{c_0}$, $b_1:=c_1$, $b_2:=c_2$.\newline Tính $\sigma:=\sqrt{\frac1n\sum_i\big(a\,e^{b_1\varphi_1(x_i)+b_2\varphi_2(x_i)}-y_i\big)^2}$.\newline Dừng thuật toán, trả về $a$, $b_1$, $b_2$, $\sigma$.
\end{thuattoan}

\begin{thuattoan}{Dạng luỹ thừa $y=a\,x^b$}
\item[Input:] vector $x$, $y$ ($n$ điểm).
\item[Output:] $a$, $b$ và sai số trung bình phương $\sigma$.
\item[Bước 0:] Đếm $n_+$ (số $y_i>0$) và $n_-$ (số $y_i<0$); đặt $s:=1$ nếu $n_+\ge n_-$, ngược lại $s:=-1$.\newline Lặp $i=1,\dots,n$:\thut nếu $s\,y_i\le0$ hoặc $x_i\le0$ thì loại điểm $i$, báo lý do.\newline Nếu còn ít hơn 2 điểm thì báo lỗi, dừng thuật toán.
\item[Bước 1:] Lặp $i=1,\dots,n$: tính $X_i:=\ln x_i$ và $z_i:=\ln(s\,y_i)$.
\item[Bước 2:] Tìm $(c_0,c_1)$ theo câu a với dữ liệu $(X_i,z_i)$ và cơ sở $(1,X)$.
\item[Bước 3:] Đặt $a:=s\,e^{c_0}$, $b:=c_1$.\newline Tính $\sigma:=\sqrt{\frac1n\sum_i(a\,x_i^b-y_i)^2}$.\newline Dừng thuật toán, trả về $a$, $b$, $\sigma$.
\end{thuattoan}

%=====================================================================
\section{Câu 23. Hàm thực nghiệm từ file dữ liệu}

\pp{Dữ liệu.} \texttt{data Câu 23.xlsx} có 4 bộ dữ liệu đặt cạnh nhau (``Câu a'' … ``Câu d''), mỗi bộ 60 cặp $(x_i,y_i)$; bộ c có mốc $x=2.63$ lặp lại với hai giá trị $y$ — giữ cả hai.

\subsection{a) $y=ax+\dfrac bx$}
\pp{Phương pháp.} Cơ sở $\varphi_1=x$, $\varphi_2=\frac1x$; hệ chuẩn:
\[ \begin{cases} a\sum x_i^2+b\sum1=\sum x_iy_i\\ a\sum1+b\sum\frac1{x_i^2}=\sum\frac{y_i}{x_i}\end{cases}\qquad(\textstyle\sum1=n=60). \]
<<c23a_bang>>
Thay các tổng:
<<c23a_he>>
Giải hệ (khử Gauss): $a=<<c23a_a>>$, $b=<<c23a_b>>$; \ $\sigma=\sqrt{\frac1{60}\sum(g(x_i)-y_i)^2}=<<c23a_sigma>>$.
\begin{ketluan}
$y\approx<<c23a_a>>\,x+\dfrac{<<c23a_b>>}{x}$, \quad $\sigma=<<c23a_sigma>>$.
\end{ketluan}

\subsection{b) $y=ax^2+b\cos x+c$}
\pp{Phương pháp.} Cơ sở $\varphi_1=x^2$, $\varphi_2=\cos x$, $\varphi_3=1$; hệ chuẩn $M_{jk}=\sum\varphi_j\varphi_k$, $b_j=\sum\varphi_jy_i$:
<<c23b_bang>>
<<c23b_he>>
Giải hệ (khử Gauss chọn trụ): $a=<<c23b_a>>$, $b=<<c23b_b>>$, $c=<<c23b_c>>$; \ $\sigma=<<c23b_sigma>>$.
\begin{ketluan}
$y\approx<<c23b_a>>\,x^2+(<<c23b_b>>)\cos x+<<c23b_c>>$, \quad $\sigma=<<c23b_sigma>>$.
\end{ketluan}

\subsection{c) $y=a\,e^{bx+cx^2}$}
\pp{Phương pháp (theo slide).} $y_i>0$ nên $z_i=\ln y_i=\ln a+bx_i+cx_i^2$, cơ sở $(1,x,x^2)$; hệ chuẩn $M_{jk}=\sum\varphi_j\varphi_k$, $b_j=\sum\varphi_jz_i$:
<<c23c_bang>>
<<c23c_he>>
Giải hệ (khử Gauss chọn trụ): $(\ln a,b,c)=<<c23c_c>>$, $a=<<c23c_a>>$. Sai số trên $z$: $\sigma_z=<<c23c_sz>>$; trên $y$ gốc: $\sigma=<<c23c_sigma>>$ — \textbf{rất lớn} so với độ lớn của $y$ ($\le2.5$).

\pp{Hiệu chỉnh.} Các điểm $x\gtrsim2.7$ có $y$ cỡ $10^{-3}$ (nhiễu) nhưng $\ln y$ sai tới vài đơn vị, chi phối hệ chuẩn. Dùng thuật toán hiệu chỉnh Gauss--Newton (Câu 22b) để cực tiểu đúng $\sum(y_i-g_i)^2$:
<<c23c_gn>>
Hội tụ sau $<<c23c_gn_lan>>$ lần: $a=<<c23c_a2>>$, $b=<<c23c_b2>>$, $c=<<c23c_cc2>>$, $\sigma=<<c23c_sigma2>>$.
\begin{ketluan}
Tuyến tính hoá (slide): $y\approx<<c23c_a>>\,e^{<<c23c_b>>x+<<c23c_cc>>x^2}$, $\sigma=<<c23c_sigma>>$ (không dùng được).\\
Sau hiệu chỉnh: $y\approx<<c23c_a2>>\,e^{<<c23c_b2>>x<<c23c_cc2>>x^2}$, \quad $\sigma=<<c23c_sigma2>>$.
\end{ketluan}

\subsection{d) $y=a\,x^b$}
\pp{Phương pháp.} $y_i<0$ ($s=-1$), $x_i>0$ nên $z_i=\ln(-y_i)=\ln(-a)+b\ln x_i$, cơ sở $(1,\ln x)$; hệ chuẩn $M_{jk}=\sum\varphi_j\varphi_k$, $b_j=\sum\varphi_jz_i$:
<<c23d_bang>>
<<c23d_he>>
Giải hệ (khử Gauss): $(\ln(-a),b)=<<c23d_c>>$, $a=-e^{\ln(-a)}=<<c23d_a>>$, $b=<<c23d_b>>$; $\sigma=<<c23d_sigma>>$ ($\sigma_z=<<c23d_sz>>$).

\pp{Hiệu chỉnh Gauss--Newton} (Câu 22b):
<<c23d_gn>>
$a=<<c23d_a2>>$, $b=<<c23d_b2>>$, $\sigma=<<c23d_sigma2>>$.
\begin{ketluan}
Tuyến tính hoá: $y\approx<<c23d_a>>\,x^{<<c23d_b>>}$, $\sigma=<<c23d_sigma>>$; \quad hiệu chỉnh: $y\approx<<c23d_a2>>\,x^{<<c23d_b2>>}$, $\sigma=<<c23d_sigma2>>$.
\end{ketluan}

\end{document}
